\subsection{COMPOSITION OF TWO FUNCTIONS. INVERSE FUNCTION}

PROPOSITION 6. If \(f\) is a mapping of A into B, and \(g\) is a mapping of B into C, then \(g \circ f\) is a mapping of A into C.

Let F, G be the graphs of f, g, respectively, and let us show that G o F is a functional graph. Let x, z, \(z'\) be objects such that \((x, z) \in \mathbf{G} \circ \mathbf{F}\) , \((x, z') \in \mathbf{G} \circ \mathbf{F}\) . There exist objects y, \(y'\) such that

\[
(x, y) \in \mathrm{F}, \quad (x, y ^ {\prime}) \in \mathrm{F}, \quad (y, z) \in \mathrm{G}, \quad (y ^ {\prime}, z ^ {\prime}) \in \mathrm{G}.
\]

Since F is a functional graph, we have \(y = y'\) and hence \((y, z') \in G\) . Since G is a functional graph, it follows that \(z = z'\) , which proves our assertion. Furthermore, the domain of \(g \circ f\) is evidently A, and the proof is complete.

The function \(g \circ f\) may also be written as \(x \to g(f(x))\) , or as \(gf\) when there is no risk of confusion.

DEFINITION 10. Let f be a mapping of A into B. The mapping f is said to be injective, or an injection, if any two distinct elements of A have distinct images under f.~The mapping f is said to be surjective, or a surjection, if \(f(A) = B\) . If f is both injective and surjective, it is said to be bijective, or a bijection.

Instead of saying that f is surjective, we sometimes say that f is a mapping of A onto B, or that it is a parametric representation of B by means of A (in which case A is called the set of parameters of the representation, and the elements of A are called parameters). If f is bijective, we sometimes say that f puts A and B in one-to-one correspondence. A bijection of A onto A is called a permutation of A.

\minorheading{Examples}

\begin{enumerate}
\def\labelenumi{(\arabic{enumi})}
\item
  If \(A \subset B\) , the mapping of \(A\) into \(B\) whose graph is the diagonal of \(A\) is injective and is called the canonical mapping or the canonical injection (or simply the injection) of \(A\) into \(B\) .
\item
  Let A be a set. The mapping \(x \to (x, x)\) of A into the diagonal \(\Delta_{A}\) of \(A \times A\) is a bijective mapping, called the diagonal mapping of A.
\item
  Let G be a set of ordered pairs. The mapping \(pr_{1}\) (resp. \(pr_{2}\) ) of G into \(pr_{1}G\) (resp. \(pr_{2}G\) ) is surjective; \(pr_{1}\) is injective if and only if G is a functional graph.
\item
  Let G be a set of ordered pairs. The mapping
\end{enumerate}

\[
z \rightarrow (\mathrm{pr} _ {2} z, \mathrm{pr} _ {1} z)
\]

of \(\mathbf{G}\) into \(\overline{\mathbf{G}}\) is a bijection (called the canonical bijection).

\begin{enumerate}
\def\labelenumi{(\arabic{enumi})}
\setcounter{enumi}{4}
\tightlist
\item
  Let \(\mathbf{A}\) be a set and \(b\) an object. The mapping \(x \to (x, b)\) of \(\mathbf{A}\) into \(\mathbf{A} \times \{b\}\) is a bijection.
\end{enumerate}

PROPOSITION 7. Let \(f\) be a mapping of A into B. Then \(\overline{f}^1\) is a function if and only if \(f\) is bijective.

If \(\overline{f}^{1}\) is a function, its source B is equal to its domain, i.e., to \(f(A)\) ; hence f is surjective. To show that f is injective, let x and y be two elements of A such that \(f(x)=f(y)\) . If F denotes the graph of f, we have

\[
(f (x), x) \in \overline {{\mathrm{F}}} ^ {- 1} \quad \text { and } \quad (f (y), y) \in \overline {{\mathrm{F}}},
\]

hence

\[
(f (x), y) \in \overline {{\mathbf {F}}},
\]

so that \(x = y\) , which proves the assertion.

Conversely, if \(f\) is bijective, it is immediate that \(\overline{\mathbf{F}}\) is functional and that the domain of \(\overline{f}^1\) is equal to \(\mathbf{B}\) .

When \(f\) is bijective, \(\overline{f}^1\) is called the inverse mapping of \(f\) ; \(\overline{f}^1\) is bijective, \(\overline{f}^1 \circ f\) is the identity mapping of A, and \(f \circ \overline{f}^1\) is the identity mapping of B.

If a permutation is the same as its inverse, it is said to be involutory.

Remark. Let \(f\) be a mapping of \(A\) into \(B\) . For each subset \(X\) of \(A\) we have (no. 3) \(X \subset f\langle f(X)\rangle\) . Furthermore, for each subset \(Y\) of \(B\) we have \(f\langle f(Y)\rangle \subset Y\) , for the relation \(y \in f\langle f(Y)\rangle\) is equivalent to

\[
(\exists x) ((\exists z) (z \in Y \text { and } z = f (x)) \text { and } y = f (x))
\]

and therefore implies the relation \((\exists z)(z \in \mathbf{Y}\) and \(y = z)\) and consequently also the relation \(y \in \mathbf{Y}\) .

If \(f\) is surjective, we have \(f\langle f\langle \mathrm{Y}\rangle\rangle = \mathrm{Y}\) for every subset \(\mathrm{Y}\) of \(\mathrm{B}\) , for the relation \(y\in \mathrm{Y}\subset \mathrm{B}\) implies by hypothesis the relation \((\exists x)(y = f(x))\) and therefore also \((\exists x)(y\in \mathrm{Y}\) and \(y = f(x)\) ); but '' \(y\in \mathrm{Y}\) and \(y = f(x)\) '' implies \((\exists z)(z\in \mathrm{Y}\) and \(z = f(x)\) ), and our assertion follows.

If \(f\) is injective, we have \(f\langle f\langle \mathbf{X}\rangle = \mathbf{X}\) for every subset \(\mathbf{X}\) of \(\mathbf{A}\) . For the relation \(x\in f^{-1}\langle f\langle \mathbf{X}\rangle\rangle\) is equivalent to \(f(x)\in f\langle \mathbf{X}\rangle\) , hence to

\[
(\exists z) (z \in \mathrm{X} \text {   and   } f (z) = f (x));
\]

but the hypothesis means that \(f(z) = f(x)\) implies \(z = x\) , hence \(x \in f\langle f\langle \mathbf{X}\rangle \rangle\) implies \(x \in \mathbf{X}\) .

